\section{Extensions and decompositions}\label{sec:extensions}%\noindent
It remains to prove the existence of extensions and
decompositions. This is the purpose of Sections~\ref{sec:extensions}
and~\ref{sec:proof}. 

First we discuss the extension problem. Let $\bB$ be a given fixed
invariant in $\E_\Bbbk(n,r-1)$. Suppose that $\bA \in \E_\Bbbk(n,r)$
is an extension of the given $\bB$, and write $\bA =
(a^\bi_\bj)_{\bi,\bj \in I(n,r)}$. Then the matrix coordinates
$a^\bi_\bj$ of $\bA$ must satisfy the conditions:
\begin{align}
 & a^{(i_1\cdots i_{r-1} i_{r-1})^\sigma}_{(j_1\cdots j_{r-1}
 j_{r-1})^\sigma} = b^{i_1\cdots i_{r-1}}_{j_1\cdots j_{r-1}} \text{ for
 all } \sigma \in \Sym_r \, , \tag{I-1} \label{I-1}\\
 & a^{i_1\cdots i_r}_{j_1\cdots j_r} = 0 \text{ whenever }
 \vt(i_1\cdots i_r) \ne \vt(j_1 \cdots j_r) \hspace{2cm}\tag{I-2}\label{I-2}
\end{align}
for all $i_1\cdots i_{r-1}$, $j_1\cdots j_{r-1}$ in $I(n,r-1)$ and all
$i_1\cdots i_r$, $j_1\cdots j_r$ in $I(n,r)$.

Condition~\eqref{I-1} comes from Proposition~\ref{prop:inv-properties}\ref{prop3.4_d},
while condition~\eqref{I-2} restates value-type preservation, from
Proposition~\ref{prop:GHS}\ref{prop3.2_b} (also Proposition~\ref{prop:inv-properties}\ref{prop3.4_c}). We call~\eqref{I-1}, \eqref{I-2}~\emph{initialisation conditions}. They determine the value of all
entries $a^\bi_\bj = a^{i_1 \cdots i_r}_{j_1\cdots j_r}$ for which
either $\#(\bi)<r$ or $\#(\bj)<r$, where we define $\#(\bi)$ to be the
number of distinct values appearing in the multi-index $\bi$. 

Thus, to find an extension $\bA$ of the given $\bB$, we start by
initialising $\bA \in \Mat_{I(n,r)}(\Bbbk)$ to satisfy~\eqref{I-1},
\eqref{I-2}. Then it only remains to assign values to the $a^{i_1 \cdots
 i_r}_{j_1 \cdots j_r}$ for which $\#(i_1 \cdots i_r) = r = \#(j_1
\cdots j_r)$. In other words, if we set
\[
I'(n,r) = \{\bi \in I(n,r): \#(\bi) = r\}
\]
then we can focus just on how to assign entries $a^\bi_\bj$ such that
$\bi,\bj \in I'(n,r)$. By Proposition~\ref{prop:GHS}, those entries of
$\bA$ must satisfy the following conditions. For any $\bp = i_1
\cdots i_{\alpha-1} i_{\alpha+1} \cdots i_r$, $\bq = j_1 \cdots
j_{\alpha-1} j_{\alpha+1} \cdots j_r$ in $I'(n,r-1)$,
\begin{align}
\textstyle \sum a^{i_1\cdots i_{\alpha-1} \, i
 \, i_{\alpha+1}\cdots i_r}_{j_1\cdots j_{\alpha-1} \,*\,
 j_{\alpha+1}\cdots j_r} = b^\bp_\bq \text{ for all } i = 1, \dots,
n\label{eq:slice-row} \\
\textstyle \sum a^{i_1\cdots i_{\alpha-1} \,*\, i_{\alpha+1}\cdots
 i_r}_{j_1\cdots j_{\alpha-1} \, j \, j_{\alpha+1}\cdots j_r} =
b^\bp_\bq \text{ for all } j = 1, \dots,
n \label{eq:slice-col}\\
\intertext{and for all $\bi = i_1 \cdots i_r$,
 $\bj = j_1 \cdots j_r$ in $I'(n,r)$,}
a^{i_1 \cdots i_r}_{j_1 \cdots j_r} - a^{(i_1 \cdots
 i_r)^\sigma}_{(j_1 \cdots j_r)^\sigma} = 0 \text{ all } \sigma \in
\Sym_r \label{eq:stability}.
\end{align}
So finding an extension $\bA$ of the given $\bB$ is (after
initialisation) equivalent to solving the linear system given by
equations~\eqref{eq:slice-row}--\eqref{eq:stability}.

Any special invariant $\bA$ in $\E_\Bbbk(n,r)^i_j$ is necessarily an
extension of its block $\bA^i_j$. This follows from Proposition~\ref{prop:special-ext-form}\ref{prop4.8_a} and Remark~\ref{rmk:block-sums}. So we
sometimes refer to such special invariants as \emph{special
 extensions}.

Our goal now is to establish the following four interrelated
properties, the first of which is about the existence of
extensions. They will be established by an interleaved double
induction on $n,r$. Note that that each property is based on some
fixed (but arbitrary) block row or column.


\begin{property}[extension property]\label{1}
 For any $\bB$ in $\E_\Bbbk(n,r-1)$, there exists some $\bA$ in
 $\E_\Bbbk(n,r)$ such that $\rho(\bA) = \bB$. More precisely, for any
 fixed $1\le i \le n$ \textup{(}respectively, $1 \le j \le
 n$\textup{)}, there exist $\bA(j)$ \textup{(}\resp
 $\bA(i)$\textup{)} in $\E_\Bbbk(n,r)^i_j$ such that $\bA = \bA(1) +
 \cdots + \bA(n)$. Call such an $\bA$ an $i$th block row
 \textup{(}\resp $j$th block column\textup{)} extension of $\bB$.
\end{property}

%\looseness-1
A priori, it is not clear that every extension can be constructed as a
block row or column extension. However, the next property guarantees
that every extension so arises.

\begin{property}[decomposition property]\label{2}
 For any given $\bA$ in $\E_\Bbbk(n,r)$ and any fixed $1\le i \le n$
 \textup{(}respectively, $1 \le j \le n$\textup{)}, there exist
 $\bA(j)$ \textup{(}\resp $\bA(i)$\textup{)} in $\E_\Bbbk(n,r)^i_j$
 such that $\bA = \bA(1) + \cdots + \bA(n)$. Call such a
 decomposition an $i$th block row \textup{(}\resp $j$th block
 column\textup{)} decomposition.
\end{property}

Property~\ref{1} for $(n,r)$ immediately implies Property~\ref{2} for
$(n,r-1)$; this follows from Remark~\ref{rmk:block-sums}, because if
$\bA$ is in $\E_\Bbbk(n,r)$ and $\rho(\bA)=\bB$ then the sum of any
chosen block row or column of $\bA$ is equal to $\bB$. It should be
noted however that we work in the opposite direction: we need
Property~\ref{2} for $(n,r-1)$ as an inductive hypothesis in order to
prove Property~\ref{1} for $(n,r)$.

Now we introduce free patterns, which index a set of entries in a
matrix which can be freely assigned to arbitary values in the ring
$\Bbbk$. Free patterns correspond to a choice of free variables in
the solution of a consistent linear system. Once the free pattern
entries have been assigned, the system has a unique solution. 


\begin{property}[free patterns for extensions]\label{3}
 For any fixed index $1\le i \le n$ \textup{(}respectively, $1 \le j
 \le n$\textup{)}, there exists a subset $F(n,r)$ of $I'(n,r) \times
 I'(n,r)$, possibly empty and depending on the index, such that for
 any $\bB$ in $\E_\Bbbk(n,r-1)$ and any assignment $f: F(n,r) \to
 \Bbbk$ there exists a unique $\bA$ in $\E_\Bbbk(n,r)$ with
 $\rho(\bA) = \bB$.
\end{property}


\begin{rema}
We often identify the elements of $F(n,r)$ with the matrix entries
they index. 
\end{rema}


It follows from the symmetry in equations
\eqref{eq:slice-row}--\eqref{eq:stability} or from the existence of
the actions in~\eqref{eq:W_n-actions} that if $F(n,r)$ is a given free
pattern, then by applying an arbitrary permutation $y\in W_n$ to all
its row or column indices, we obtain another free
pattern. Furthermore, by interchanging row and column indices in
$F(n,r)$ (transposing) we obtain another free pattern.

\begin{property}[free patterns for decompositions]\label{4}
 For any given $\bA$ in $\E_\Bbbk(n,r)$ and any fixed $1\le i \le n$
 \textup{(}respectively, $1 \le j \le n$\textup{)}, there exists a
 subset $D(n,r)$ of $\{1, \dots, n\} \times I'(n,r) \times I'(n,r)$
 such that for any given assignment $f: D(n,r) \to \Bbbk$ there exist
 unique $\bA(j)$ \textup{(}\resp $\bA(i)$\textup{)} in
 $\E_\Bbbk(n,r)^i_j$ such that $\bA = \bA(1) + \cdots + \bA(n)$.
\end{property}

Properties~\ref{3},~\ref{4} are refinements of
Properties~\ref{1},~\ref{2} (respectively) that precisely quantify the
amount of freedom in constructing solutions to their underlying linear
systems.

Suppose that $\bA$ is in $\E_\Bbbk(n,r)$. Fix $i$ and consider its
$i$th block row $\bA^i_* = (\bA^i_j)_{j=1}^n$. By Lemma
\ref{lem:blocks-are-special}, each block $\bA^i_j$ of $\bA^i_*$ gives
an invariant $\bB(j)$ in $\E_\Bbbk(n,r-1)^i_j$, and the correspondence
is given by
\begin{equation}\label{eq:label-iso}
 a^{i\, i_2 \cdots i_r}_{j\,j_2 \cdots j_r} = b(j)^{i_2 \cdots i_r}_{j_2 \cdots j_r}
\end{equation}
as $j$ runs from $1$ to $n$. We let $\pi^i$ be the bijection between
the set of pairs $(i\, i_2 \cdots i_r, j\,j_2 \cdots j_r)$ indexing
entries of $\bA^i_*$ on the left hand side of~\eqref{eq:label-iso} and
the set of triples $(j, i_2 \cdots i_r, j_2 \cdots j_r)$ indexing
entries on the right hand side. Similarly, transposing rows and
columns and $i,j$ we obtain a bijection $\pi_j$ between the set of
pairs indexing entries of the $j$th block column $\bA^*_j$ and a
corresponding set of triples. The following lemma is immediate from
Remark~\ref{rmk:block-sums}.


\begin{lemm}\label{lem:bijections}
Let $\bA \in \E_\Bbbk(n,r)$ and set $\bB = \rho(\bA) \in
\E_\Bbbk(n,r-1)$. The map $\pi^i$ \textup{(}\resp $\pi_j$\textup{)}
is a bijection between a set of labels for the entries of $\bA^i_*$
\textup{(}\resp $\bA^*_j$\textup{)} and a set of labels for a solution
$\bB = \bB(1)+ \cdots + \bB(n)$ to the decomposition problem in degree
$r-1$.
\end{lemm}





Now suppose that $F(n,r)$ exists, in other words, that the free
extension pattern for constructing $\bA$ (from $\bB$) in Lemma~\ref{lem:bijections} exists (note that $F(n,r)$ is necessarily based on a designated block row or column). If $F(n,r)$ is based on the
$i$th block row (\resp $j$th block column) then we define 
\begin{equation}\label{eq:F'-def}
 F'(n,r) = \{(i_1 \cdots i_r, j_1 \cdots j_r) \in F(n,r):
 i_1 = i\ (\text{\resp}\ j_1 = j)\}.
\end{equation}
Furthermore, we define $F''(n,r)=F(n,r) \setminus F'(n,r)$, so that
\begin{equation}\label{eq:crux}
F(n,r) = F'(n,r) \sqcup F''(n,r).
\end{equation}
This disjoint decomposition is the crux of the interleaved induction
that will prove Properties~\ref{1}--\ref{4}. Of great importance for
our results is the fact that the map $\pi^i$ (\resp $\pi_j$) in Lemma~\ref{lem:bijections} restricts to a bijection
\[
F'(n,r)\cong D(n,r-1),
\]
where the right-hand side is the free decomposition pattern for writing
$\bB = \bB(1)+ \cdots + \bB(n)$ in Lemma~\ref{lem:bijections}. In
particular, this means that Property~\ref{3} for $(n,r)$ immediately
implies Property~\ref{4} for $(n,r-1)$. However, our induction
proceeds in the reverse direction, using the inverse of the above
bijection to construct $F'(n,r)$ from $D(n,r-1)$. Then we construct
$F''(n,r)$ and glue them together according to~\eqref{eq:crux} in
order to obtain $F(n,r)$. In this manner, we explicitly construct and
interrelate the free extension patterns and free decomposition
patterns.


\begin{defi}
We order $I(n,r)$ lexicographically, and we do the same
for the row and column indices of any matrix in
$\Mat_{I(n,r)}(\Bbbk)$. A free pattern $F(n,r)$ is \emph{row-initial}
(\resp \emph{column-initial}) if, after identifying free-pattern
elements with their corresponding matrix entries, free-pattern entries
precede all other entries in each row (\resp column)
slice. Similarly, it is \emph{row-terminal} (\resp
\emph{column-terminal}) if, after identifying free-pattern elements
with matrix entries, free-pattern entries come after all other entries
in each row (\resp column) slice.
\end{defi}

It is evident that free patterns $F(n,1)$ exist. Indeed, it is easily
checked that
\[
F(n,1) = \{(i,j): 2 \le i,j \le n\}
\]
is a row- and column-terminal pattern for $(n,1)$. To see this,
notice that once an assignment $f:F(n,1) \to \Bbbk$ has been chosen,
there is a unique extension $\bA$ in $\E_\Bbbk(n,1)$ of any given $b
\in \E_\Bbbk(n,0) = \Bbbk$ satisfying $a^i_j = f(i,j)$ for all $(i,j)
\in F(n,1)$. Note that $a^1_i$ and $a^i_1$ are then uniquely forced
for $ i>1 $. The first entry $a^1_1$ is forced in two ways, but
easily seen to be well-defined.

By applying appropriate permutations to the rows and/or columns of
$F(n,1)$ above, one gets row- and column-terminal, row-initial and
column-terminal, and row-terminal and column-initial free patterns for
$(n,1)$. Our proof of the four basic properties in the next section
will show that these variations of free patterns always exist, for any
$(n,r)$. Note that the distinguished element $w_0 \in W_n$ given by
$w_0(j) = n+1-j$ for $j = 1, \dots, n$ interchanges initial and
terminal patterns.


We conclude this section with the following algorithm, which (as we
show in Corollary~\ref{cor:colouring}) determines a terminal
(respectively, initial) free pattern in a randomly chosen row or
column of an extension $\bA$ in $\E_\Bbbk(n,r)$ of a given $\bB$ in
$\E_\Bbbk(n,r-1)$, assuming that it is the first such row or column to
be completed. The algorithm is independent of
Properties~\ref{1}--\ref{4}, and independent of the inductive proof of
those properties.

\begin{algorithm}\label{algo}
The $\alpha$-slices of $I'(n,r)=\{\bi \in I(n,r): \#(\bi)=r\}$ are the
subsets
\[
\{i_1 \cdots i_{\alpha-1}\, i\, i_{\alpha+1}\cdots i_r: i = 1, \dots,
n\}
\]
for $\alpha = 1, \dots, r$. We start by listing the elements of
$I'(n,r)$ in lexicographical order. We recursively assign a colour
\textup{(}$0$ or $1$\textup{)} to each element of $I'(n,r)$ as follows.

As long as uncoloured elements exist, we find the largest
\textup{(}respectively, smallest\textup{)} uncoloured element, colour
it, and repeat. Colouring an element consists of the following two
steps:
\begin{enumerate}[label=(\alph*)]
\item\label{algo5.4_i} Examine all the element's slices. If the element is the only
 uncoloured element in one of its slices, we say it is forced, and
 colour it $0$. Otherwise, colour the element $1$ to indicate that it
 is free.
\item\label{algo5.4_ii} If colouring the current element forces any additional elements
 in any of its slices, then colour those elements as
 well. \textup{(}This happens when there is just one remaining
 uncoloured element in the slice, after the current element is
 coloured.\textup{)}
\end{enumerate}
We note that colouring is a recursive process, because of~\ref{algo5.4_ii}. Let $I'_1(n,r)$ be the set of elements of $I'(n,r)$
coloured $1$ by the above procedure.
\end{algorithm}

\begin{exam}
By always choosing the largest uncoloured element, for the case of
$I'_1(5,2)$ the algorithm produces the following colouring:
\[
 I'_1(5,2): \quad
\text{\setlength{\arraycolsep}{2.7pt}\scriptsize
$\displaystyle \begin{array}{c|cccc|cccc|cccc|cccc|cccc|}
 &\rt{12}&\rt{13}&\rt{14}&\rt{15}&\rt{21} &\rt{23}&\rt{24}&\rt{25}
 &\rt{31}&\rt{32} &\rt{34}&\rt{35}&\rt{41}&\rt{42}&\rt{43} &\rt{45} 
 &\rt{51}&\rt{52}&\rt{53} &\rt{54} \\ \hline
 &0&0&0&0	&0&0&1&1	 &0&1&1&1	&0&1&1&1 &0&1&1&1
 \\ 
 \hline
\end{array}\; .$}
\]
Thus we have $I'_1(5,2) = \{54, 53, 52, 45, 43, 42, 35, 34, 32, 25, 24\}$.
\end{exam}

Let $\bB$ in $\E_\Bbbk(n,r-1)$ be given, and fix some index $\bi$
(respectively, $\bj$) in $I'(n,r)$. Then any assignment
 \[
 f: \{ a^\bi_\bj : \bj \in I'_1(n,r)\} \to \Bbbk \qquad (\resp f: \{
 a^\bi_\bj : \bi \in I'_1(n,r)\} \to \Bbbk )
 \]
determines the $\bi$th row $a^\bi_*$ (\resp $\bj$th column $a^*_\bj$)
of a matrix $\bA$ satisfying equations~\eqref{eq:slice-row} (\resp \eqref{eq:slice-col}) with respect to the given $\bB$.


\goodbreak
\begin{rema}\label{rmk:promise}\ \\*[-1.2em]
\begin{enumerate}[label=(\roman*)]
\item\label{rema5.6_i} In Corollary~\ref{cor:colouring} we will show that that, under
suitable inductive hypotheses, there exists an extension $\bA$ of the
given $\bB$ which agrees with the row or column determined by the
above algorithm.

\item\label{rema5.6_ii} We prefer to work with terminal free patterns, because they are
compatible with restriction, in the sense that by excising all indices
containing an $n$ we obtain a free pattern for $n-1$. This preference
pervades all of the examples and some of the proofs in the next
section.
\end{enumerate}
\end{rema}



